\paragraph{Assumptions.}
Throughout this section we assume that:

\begin{itemize}
\item $S$ is finite.
\item Every right-hand side occurring in $S$ is a finite type expression.
\item Right-hand sides of type definitions contain no type variables.
\item For every type variable $v$, the expression $\sigma(v)$ is variable-free.
\end{itemize}

\begin{theorem}
For all type expressions $e_1,e_2$, type mappings $\sigma$, and guard sets $C$,
the computation $\meet_S(e_1,e_2,\sigma,C)$ terminates.
\end{theorem}

\begin{proof}

Let $K$ be the set of type constants occurring in $S$. Since $S$ is finite, so is $K$.
Since $S$ is finite and all type expressions occurring in $S$, the input, and the values $\sigma(v)$ are finite, the set of their subexpressions is finite.
Every recursive call of $\meet_S$ either unfolds a type definition, replaces a variable by some $\sigma(v)$, or recurses on subexpressions of its arguments.
Hence every expression occurring as an argument of a recursive call is a subexpression of either the input, a right-hand side in $S$, or a value $\sigma(v)$.
Consequently, only finitely many expressions can occur as arguments of recursive calls.
Since both $K$ and the set of possible recursive-call arguments are finite, only finitely many guard pairs $(t,e)$ can arise during the computation. Let $N$ denote this number.
We define

\[
g(C)=N-|C|.
\]

Furthermore, let the size of a type expression be given by
\[
|\top|=|t|=|v|=1,\qquad
|f(e_1,\ldots,e_n)|=1+\sum_{i=1}^{n}|e_i|,\qquad
|e_1|\cdots|e_n|=1+\sum_{i=1}^{n}|e_i|.
\]

For a pair of expressions, let
\[
s(e_1,e_2)=|e_1|+|e_2|.
\]

We use the lexicographic measure
\[
M(e_1,e_2,C)=\bigl(g(C),s(e_1,e_2)\bigr).
\]

Since $g(C)$ and $s(e_1,e_2)$ are natural numbers, the lexicographic order on $\mathbb N\times\mathbb N$ is well-founded.

We now inspect the recursive cases of $\meet_S$.

\begin{itemize}

\item \textbf{Type unfolding.}
A recursive call of the form
\[
\meet_S(t,e,\sigma,C)\rightarrow\meet_S(e',e,\sigma,C\cup\{(t,e)\})
\]
with $t=e'\in S$ and $(t,e)\notin C$ adds a fresh guard pair to $C$. Therefore the first component of $M$ strictly decreases.

\item \textbf{Subtype unfolding.}
A recursive call of the form
\[
\meet_S(f(e_1,\ldots,e_n),t,\sigma,C)\rightarrow
\meet_S(f(e_1,\ldots,e_n),e,\sigma,C\cup\{(t,f(e_1,\ldots,e_n))\})
\]
with $t\sqsubseteq e\in S$ and $(t,f(e_1,\ldots,e_n))\notin C$ also adds a fresh guard pair and therefore
strictly decreases the first component of $M$.

\item \textbf{Variable expansion.}
A recursive call of the form
\[
\meet_S(v,e,\sigma,C)\rightarrow\meet_S(\sigma(v),e,\sigma,C)
\]
replaces the variable $v$ by its value. Since $\sigma(v)$ is variable-free, the recursive call cannot itself match
the variable-expansion case. Thus every variable expansion is immediately followed by one of the remaining cases.

\item \textbf{Union decomposition.}
For a call of the form $\meet_S(e_1|\cdots|e_n,e,\sigma,C)$, recursive calls are performed on
the pairs $(e_i,e)$. Since every $e_i$ is a strict subexpression of $e_1|\cdots|e_n$,
we have $s(e_i,e)<s(e_1|\cdots|e_n,e)$. Thus the second component of $M$ strictly decreases while the first component remains unchanged.

\item \textbf{Function decomposition.}
For a call of the form $\meet_S(f(e_1,\ldots,e_n),f(g_1,\ldots,g_n),\sigma,C)$, recursive calls are performed
on the pairs $(e_i,g_i)$. Since $e_i$ and $g_i$ are strict subexpressions of the original function types, we have
\[
s(e_i,g_i)<s(f(e_1,\ldots,e_n),f(g_1,\ldots,g_n)).
\]
Hence the second component of $M$ strictly decreases while the first component remains unchanged.

\end{itemize}

All remaining cases perform no recursive calls and therefore terminate immediately.

Thus every recursive call either strictly decreases the first component of $M$, strictly decreases the second component
while leaving the first unchanged, or is a variable expansion that immediately reduces to one of the previous cases.
Consequently, every infinite computation would yield an infinite descending chain in the well-founded order on
$\mathbb N\times\mathbb N$, which is impossible. Therefore $\meet_S(e_1,e_2,\sigma,C)$ terminates.

\end{proof}
